\uuid{7Eo5}
\titre{Maximum}
\auteur{bodin ; hanani ; mzari}
\organisation{exo7}
\licence{CC BY-NC-SA 4.0}
\datecreate{2020-01-20}
\niveau{L1}
\module{Analyse}
\chapitre{Propriétés de R}
\sousChapitre{Propriétés des nombres réels}
\difficulte{2}
\theme{Propriétés de nombres réels, maximum}
\exo7qcmid{lille-1.reels.011}

\contenu{
\begin{question}
\qclassification{120.03}
\qtags{motcle=maximum}
Quelles sont les assertions vraies concernant le maximum de nombres réels ?
\begin{answers}
    \good{$\max(a,b) \ge a$ et $\max(a,b) \ge b$}

    \bad{$\max(a,b) > a$ ou $\max(a,b) > b$}

    \good{$\max( \max(a,b), c ) = \max(a,b,c)$}

    \good{$\min(a, \max(a,b)) = a$}
\end{answers}
\begin{explanations}
$\max(a,b) \ge a$ et $\max(a,b) \ge b$ et $\max( \max(a,b), c ) = \max(a,b,c)$. L'assertion "$\max(a,b) > a$ ou $\max(a,b) > b$" est fausse (prendre $a=b$). Cherchez une preuve pour l'assertion restante !
\end{explanations}
\end{question}
}
